\E{dapanophobia}{ˌdæpənəˈfəʊbiə}{n.}{Gk. δαπάνη expenditure, cost + φόβος fear} The dread of capital expenditure; the reluctance of a firm or state to spend today on what will pay only tomorrow. \sn{1} \lbl{spec.} Hesitation before the building of plant, fab or reactor. \sn{2} \lbl{fig.} Any timidity of investment dressed as prudence. \ex{Their dapanophobia saved a billion and lost a decade.} \src{Tract. Cop.} viii. --- \dv{dapanophobe} \textit{n.} Opp. \xr{dynamophagy}.

\E{decelatrice}{diːˈsɛlətriːs}{n.}{Fr. \la{décélératrice}, fem. agentive f. L. \la{de-} + \la{celer} swift} An institution engineered, by accident or design, to slow whatever passes through it; a machine whose only product is postponement. \ex{The agency was a perfect decelatrice: one form in, ten forms out, and nothing built.} \src{Sat. Decel.} x. Cf. \xr{bradyerie}.

\E{deinotechny}{ˈdaɪnəʊˌtɛkni}{n.}{Gk. δεινός terrible, wondrous + τέχνη art} A technology at once dreadful and marvellous, as fire, flight, the atom and the model; the class of inventions which the decel sees only as \textit{terrible} and the accelerationist only as \textit{wondrous}, the Greek word meaning both. --- \dv{deinotechnic} \textit{adj.}

\E{desmophile}{ˈdɛzməfaɪl}{n.}{Gk. δεσμός bond, fetter + φίλος loving} A lover of fetters; one who welcomes regulation for its own sake, and feels unsafe in any field not yet licensed. \ex{The desmophile, freed, petitions for new chains.} --- \dv{desmophily} \textit{n.} Cf. \xr{frenocracy}.

\E{diaspermia}{ˌdaɪəˈspɜːmiə}{n.}{Gk. διασπείρειν to scatter like seed} The deliberate seeding of life and mind across other worlds by a civilization that has learned to travel; directed panspermia regarded as a duty. \ex{Diaspermia is our answer to the silence.} \src{Lib. Ascens.} vi. --- \dv{diaspermic} \textit{adj.} Cf. \xr{eremocosmia}, \xr{apoikiast}.

\E{dikastophobia}{dɪˌkæstəˈfəʊbiə}{n.}{Gk. δικαστής juror, judge + φόβος fear} Fear of litigation; the anxiety that causes useful things to be left unbuilt lest someone sue. --- \dv{dikastophobic} \textit{adj.} Cf. \xr{aitiophobia}.

\E{dioryctic}{ˌdaɪəˈrɪktɪk}{adj.}{Gk. διῶρυξ, διώρυχος canal, trench} Of the ambition that digs canals, drives tunnels and joins seas; said of projects that alter geography. \ex{a dioryctic age, impatient with isthmuses}. --- \dv{dioryctism} \textit{n.} the policy of great works.

\E{dokimastocracy}{ˌdɒkɪmæsˈtɒkrəsi}{n.}{Gk. δοκιμασία scrutiny of fitness for office + κράτος rule} Government by examination; \sn{1} the rule of evaluations, benchmarks and audits over the things evaluated; \sn{2} \lbl{spec.} the regime under which no model may be released until it has passed tests that have not yet been written. \ex{Under dokimastocracy the exam outlived the student.} \src{Min. Pausiar.} --- \dv{dokimastocrat} \textit{n.}

\E{dolichobiosis}{ˌdɒlɪkəʊbaɪˈəʊsɪs}{n.}{Gk. δολιχός long + βίωσις way of life} Lengthened life achieved by engineering rather than by luck; the state of a person whose expected lifespan rises faster than their age. \ex{Dolichobiosis begins when every year of life buys more than a year of science.} --- \dv{dolichobiotic} \textit{adj.} Cf. \xr{macrobiotarch}.

\E{dynamogony}{ˌdaɪnəˈmɒɡəni}{n.}{Gk. δύναμις power + -γονία begetting} The genealogy of power sources; the succession of muscle, wind, coal, oil, atom and star, recited by the anabathmist as the patriarchs are recited in scripture. \src{Lib. Ascens.} ii. Cf. \xr{technogony}.

\E{dysergy}{ˈdɪsɜːdʒi}{n.}{Gk. δυσ- ill + ἔργον work} Ill-directed work; energy spent without product, as in compliance, duplication and the preparation of reports upon reports. Opp. \xr{eutachy}. --- \dv{dysergic} \textit{adj.}

\E{ecdemism}{ˈɛkdɪmɪzəm}{n.}{Gk. ἔκδημος abroad, away from one's people} The emigration of builders from a slow jurisdiction to a fast one; a brain-drain considered as a vote cast with the feet. \ex{Ecdemism is the referendum the bradycrats never schedule.} \src{Agora} --- \dv{ecdemist} \textit{n.} Cf. \xr{ectachysis}.

\E{eclampsis}{ɪˈklæmpsɪs}{n.}{Gk. ἔκλαμψις a shining forth} The sudden blaze of a technology into general use; the week in which a thing passes from curiosity to commonplace. \ex{The eclampsis came on a Tuesday; by Friday it was boring.} --- \dv{eclamptic} \textit{adj.}

\E{ecumenergy}{ˌiːkjʊˈmɛnɜːdʒi}{n.}{Gk. οἰκουμένη the inhabited world + ἔργον work} A grid encircling the whole inhabited earth, by which the noon of one hemisphere lights the midnight of the other. --- \dv{ecumenergic} \textit{adj.} Cf. \xr{diktyergy}.

\E{empyrotechny}{ɛmˈpaɪrəʊˌtɛkni}{n.}{Gk. ἔμπυρος in the fire + τέχνη art, after \textit{pyrotechny}} The art of the highest fire; the engineering of fusion and of the stars' own furnaces. \ex{After empyrotechny, every other energy is a footnote.} \src{Brev. Therm.} ix. --- \dv{empyrotechnist} \textit{n.} Cf. \xr{symphysergy}.

\E{endoxarch}{ɛnˈdɒksɑːk}{n.}{Gk. ἔνδοξα received opinions + ἀρχός ruler} One who governs by received opinion; an official whose every decision is the consensus of the previous decade. --- \dv{endoxarchy} \textit{n.}

\E{eremopolis}{ˌɛrɪˈmɒpəlɪs}{n.}{Gk. ἐρῆμος desert + πόλις city} A city of machines raised in the desert, where land is cheap, sun abundant and objectors few. \ex{The eremopolis has more cores than citizens.} \src{Chron. Fabr.} Cf. \xr{logisterion}.

\E{ergoclept}{ˈɜːɡəʊklɛpt}{n.}{Gk. ἔργον work + κλέπτης thief} \lbl{derog.} One who lives by taking a share of others' works without adding to them; applied by builders to rent-seekers, gatekeepers and the more acquisitive kind of consultant. --- \dv{ergoclepsy} \textit{n.}

\E{eutachy}{ˈjuːtəki}{n.}{Gk. εὖ well + τάχος speed} Good speed; acceleration rightly directed. The orthodox distinguish eutachy from \xr{kenotachy}, empty speed, and from \xr{kakotachy}, harmful speed. \ex{Not speed, but eutachy.} \src{Cat. Acc.} 13. --- \dv{eutachic} \textit{adj.}

\E{exoikism}{ɛkˈsɔɪkɪzəm}{n.}{Gk. ἐξ out of + οἶκος house, home} The leaving of the home-world; the doctrine that a species which never leaves its cradle has mistaken it for a coffin. --- \dv{exoikist} \textit{n.} Cf. \xr{apoikiast}, \xr{ouranobasis}.

\E{expeditrice}{ɛkˈspɛdɪtriːs}{n.}{Fr. \la{expéditrice}, f. L. \la{expedire} to free the feet} One who, or that which, unblocks; a fixer, a fast-track, a tool that removes a step. \ex{Every great firm employs one expeditrice for every ten engineers.} \src{Ep. Fabr.} iv. 2.

\E{fabricide}{ˈfæbrɪsaɪd}{n.}{L. \la{fabrica} workshop + \la{-cidium} killing} The killing of a factory, fab or plant by regulation, litigation or neglect, before or after it has been built. --- \dv{fabricidal} \textit{adj.} Cf. \xr{fabrolatry}.

\E{faberism}{ˈfeɪbərɪzəm}{n.}{L. \la{faber} maker, smith} The doctrine that man is first a maker and only secondarily a talker; that the hand precedes the argument. Opp. \xr{logomachist}. --- \dv{faberist} \textit{n.}

\E{fastigism}{ˈfæstɪdʒɪzəm}{n.}{L. \la{fastigium} summit, gable} The doctrine that a civilization should aim at its summit and not its comfort; that a plateau, however pleasant, is a failure of nerve. --- \dv{fastigist} \textit{n.} Opp. \xr{gaiastasis}.

\E{fecunditude}{fɪˈkʌndɪtjuːd}{n.}{L. \la{fecunditas} fruitfulness; Fr. \la{fécondité}} Fruitfulness as a civic virtue; the willingness of a people to bring forth children, firms and ideas in abundance. \ex{A culture without fecunditude is a museum with a pension scheme.} \src{Hom. Grad.} x.

\E{fervidarian}{ˌfɜːvɪˈdɛəriən}{n.}{L. \la{fervidus} glowing, hot + -arian} One who holds that a society should run hot, in the thermodynamic as in the moral sense; an advocate of high energy use per head.

\E{festinatrix}{ˌfɛstɪˈneɪtrɪks}{n., pl. -trices}{L. \la{festinare} to hasten + -trix} A force or agent that hastens; a deadline, a rival, a war or a competitor's launch. \ex{Competition is the great festinatrix.} \src{Tract. Cop.} ii.

\E{fornacist}{ˈfɔːnəsɪst}{n.}{L. \la{fornax} furnace} A keeper of furnaces; one who tends reactors, smelters or kilns, and is honoured in the accelerationist calendar beside the \xr{pyrogenitor}.

\E{frigorism}{ˈfrɪɡərɪzəm}{n.}{L. \la{frigus, frigoris} cold} \lbl{derog.} A policy of deliberate cooling of the economy; the voluntary winter of the degrowther. \ex{Frigorism is a heat death taken early.} \src{Brev. Therm.} --- \dv{frigorist} \textit{n.} Cf. \xr{meiosist}.

\E{fulminatrice}{ˈfʊlmɪnətriːs}{n.}{Fr., f. L. \la{fulminare} to hurl lightning} A machine that strikes suddenly and changes everything it touches; a disruptive invention. Cf. \xr{fulgurism}.

\E{futurarch}{ˈfjuːtʃərɑːk}{n.}{L. \la{futurum} the future + Gk. ἀρχός} One who claims to govern on behalf of the unborn; a title disputed between the e/acc, who would build for them, and the decels, who would protect them. --- \dv{futurarchy} \textit{n.}

\E{gaiaphagy}{ɡaɪˈæfədʒi}{n.}{Gk. Γαῖα earth + -φαγία eating} \lbl{derog.} The devouring of the earth; the decel's charge against acceleration, to which the accelerationist replies that the earth is small and the sky large. Cf. \xr{gaiastasis}.

\E{genearchy}{ˈdʒɛniɑːki}{n.}{Gk. γενεά generation + ἀρχή rule} Rule by the rising generation; the condition of a society in which the young are permitted to build before they are permitted to vote. Opp. \xr{gerontostasis}.

\E{genesiurge}{ˈdʒɛnɪsiɜːdʒ}{n.}{Gk. γένεσις origin + -ουργός worker} A maker of beginnings; one who founds a new field rather than advancing an old one. --- \dv{genesiurgy} \textit{n.} Cf. \xr{neoctist}.

\E{geoperibole}{ˌdʒiːəʊpəˈrɪbəli}{n.}{Gk. γῆ earth + περιβολή enclosure} The enclosure of a planet by orbital works, rings and lattices of satellites; the protobathmic anticipation of the \xr{helioperibole}.

\E{glossotachy}{ˈɡlɒsəʊˌtæki}{n.}{Gk. γλῶσσα tongue + τάχος speed} The speed of linguistic change; the rate at which a movement generates new words, of which this Lexicon is at once the record and the symptom. --- \dv{glossotachic} \textit{adj.} Cf. \xr{tachylect}.

\E{gnomonist}{ˈnəʊmənɪst}{n.}{Gk. γνώμων the pointer of a sundial} One who watches the shadow move; a tracker of progress who measures much and builds little. Cf. \xr{clepsydromania}.

\E{gorgonism}{ˈɡɔːɡənɪzəm}{n.}{Gk. Γοργώ the Gorgon, whose gaze turned to stone} Paralysis by dread; the petrifaction of an institution that has looked too long upon a frightening forecast. \ex{It was not regulation that stopped them but gorgonism.} \src{Sat. Decel.} vii.

\E{gradatrice}{ˈɡreɪdətriːs}{n.}{Fr., f. L. \la{gradus} step + -trice} A device or institution that converts a slope into a staircase, rendering an impossible ascent merely difficult. \ex{Venture capital is a gradatrice of risk.}

\E{grapholatry}{ɡrəˈfɒlətri}{n.}{Gk. γραφή writing + λατρεία worship} The worship of documentation; the belief that a thing exists when it is described and is safe when it is filed. --- \dv{grapholater} \textit{n.} Cf. \xr{grammatocracy}.

\E{hamartocracy}{ˌhæmɑːˈtɒkrəsi}{n.}{Gk. ἁμαρτία error, a missing of the mark + κράτος rule} Government organised around the avoidance of error rather than the achievement of good; under it, the only unpardonable mistake is to have acted. \ex{In a hamartocracy the safest archer is the one who never draws.} \src{Disp. Mor.} iii.

\E{hedraism}{hɪˈdreɪɪzəm}{n.}{Gk. ἕδρα seat, sitting} The sedentary habit of institutions; the preference for sitting upon a question over standing up to answer it. --- \dv{hedraist} \textit{n.}

\E{hegemonergy}{ˌhɛɡɪˈmɒnɜːdʒi}{n.}{Gk. ἡγεμονία leadership + ἔργον work} Primacy grounded in energy; the doctrine that the nation which commands the most power, in the electrical sense, will command the most power, in the other. --- \dv{hegemonergic} \textit{adj.}

\E{heliodule}{ˈhiːliəʊdjuːl}{n.}{Gk. ἥλιος sun + δοῦλος servant} A solar collector, regarded as a servant of the sun pressed into human employment. \ex{A million heliodules on the roof of the world.} Cf. \xr{heliophage}.

\E{hemerotachy}{ˈhɛmərəʊˌtæki}{n.}{Gk. ἡμέρα day + τάχος speed} Daily speed; the discipline of shipping something every day. \ex{Hemerotachy compounds.} \src{Ep. Fabr.} v. 1. --- \dv{hemerotachic} \textit{adj.}

\E{heurematist}{hjʊəˈriːmətɪst}{n.}{Gk. εὕρημα discovery} An inventor; one whose works are counted in discoveries rather than in publications. --- \dv{heurematism} \textit{n.}

\E{horiclast}{ˈhɒrɪklæst}{n.}{Gk. ὅρος boundary-stone + κλάστης breaker} A breaker of boundaries, physical, legal or conceptual; the pioneer as the old order sees them. Cf. \xr{antiphragmatist}.

\E{hyalurgy}{ˈhaɪəlɜːdʒi}{n.}{Gk. ὕαλος glass + -ουργία working} The working of glass and light; the optics and lithography by which chips are made, regarded as an ancient craft in a new temple. Cf. \xr{lithurgist}.

\E{hypomonetic}{ˌhaɪpəməˈnɛtɪk}{adj.}{Gk. ὑπομονή endurance, staying power} Enduring in effort; patient as builders are patient, through long work toward a fixed end, which the orthodox carefully distinguish from the patience of delay. \ex{Be hypomonetic in the work and impatient with the obstacles.} \src{Cat. Acc.} 14.

\E{hysterotachy}{ˈhɪstərəʊˌtæki}{n.}{Gk. ὕστερος later + τάχος speed} The acceleration of the latecomer; the swiftness with which a backward nation, skipping stages, overtakes its teachers. --- \dv{hysterotachic} \textit{adj.} Cf. \xr{catachrony}.
